\documentclass[11pt]{article}
\usepackage{fontspec}
\setmainfont{Times New Roman}
\setsansfont{Arial}
\setmonofont{Consolas}
\usepackage{amsmath,amssymb,mathtools}
\usepackage{geometry}
\usepackage{microtype}
\geometry{a4paper,margin=2.28cm}
\setlength{\parindent}{1.15em}
\setlength{\parskip}{0pt}
\makeatletter
\renewcommand\footnotesize{\@setfontsize\footnotesize{7.5}{8.5}\selectfont}
\makeatother
\providecommand{\ma}{\mathfrak{a}}
\providecommand{\mc}{\mathfrak{c}}
\providecommand{\md}{\mathfrak{d}}
\providecommand{\ml}{\mathfrak{l}}
\providecommand{\mo}{\mathfrak{o}}
\providecommand{\mr}{\mathfrak{r}}
\providecommand{\mt}{\mathfrak{t}}
\emergencystretch=120em
\allowdisplaybreaks
\sloppy
\hbadness=10000
\vbadness=10000
\hfuzz=2pt
\vfuzz=10pt
\raggedbottom
\begin{document}
% Unit: Paper34 section13 v001. Direct Interslavic Latin translation from audited German source lines 772--882, segments P34-S0150--P34-S0166. Source scan pages 23--25 retained as witnesses; Section14 heading P34-S0167 is intentionally reserved for Section14. Zenodo record 20822156 was rechecked for this unit at 2026-06-24T23:34Z UTC with no file key/size/checksum change.
\section*{34. Hiperkompleksne veličiny i teorija predstavjenj}
\emph{Prodolženje: glava II, §§8--14.}

\section*{II. poglavje. Nekomutativna teorija idealov}

\subsection*{§ 13. Popolno reducibilne kolca}

Po §5 kolco nazyva se pravo popolno reducibilnym, ako ono jest direktna suma konečno mnogih prostyh pravyh idealov.

\emph{Pravo popolno reducibilno kolco s jediničnym elementom ne imaje nilpotentnogo ideala $\ne(0)$, zato ne imaje radikala.}

\emph{Dokaz.} Vsaky pravy ideal $\mr$ jest direktny sumand (§5):
\[
        \mo=\mt+\mr=e_1\mo+e_2\mo.
\]
Poneže $e_2^2=e_2$, imajemo takože $e_2^\rho=e_2$. Ako by $\mr$ byl nilpotentny, togda bylo by $e_2^\rho=0$, zato $e_2=0$, zato $\mr=0$, q.e.d.

Teper pokažemo oba obratna tvrđenja. Prvo: \emph{kolco bez radikala, s uslovjem minimalnosti za prave idealy, jest popolno reducibilno v odnosu do pravyh idealov.}

\emph{Dokaz.} Nehaj $\mt$ bude minimalny pravy ideal $\ne0$. Hočemo pokazati, že $\mt$ sadrži idempotentny element.

Poneže $\mt^2\subseteq\mt$ i $\mt^2\ne0$, slijedi $\mt^2=\mt$, bo $\mt^2$ jest pravy ideal. Zato jest nekoje $a$ v $\mt$ tako, že $a\mt\ne0$. Tada mora byti $a\mt=\mt$, bo $a\mt$ jest pravy ideal.

Vse elementy $b$ v $\mt$, ktore anulirajut se črez $a$ ($ab=0$), tvorjat pravy ideal, i ne jest on $=\mt$, zato jest $=0$. Znači: iz $ab=0$ slijedi $b=0$.

Zaradi $\mt=a\mt$ element $a$ možno predstaviti v obliku $ac$: $a=ac$, $c\ne0$. Slijedi
\[
        ac=ac^2,
        \qquad a(c-c^2)=0,
        \qquad c-c^2=0,
        \qquad c^2=c.
\]
Teper dalje pokažemo, že $\mt$ jest direktny sumand. $c\mo$ jest pravy ideal $\subseteq\mt$, $\ne0$, bo $c^2=c$ leži v $c\mo$; zato $c\mo=\mt$. Vsaky element $r$ iz $\mo$ dopušča predstavjenje
\[
        r=cr+(r-cr) \qquad \hbox{(jednostranno Peirceovo razloženje).}
\]
Elementy $cr$ tvorjat ideal $\mt$; elementy $r-cr$ tvorjat drugy pravy ideal $\mr$, ktory jest anulovany črez $c$: $c\mr=0$. Zato
\[
        \mo=(\mt,\mr).
\]
Poneže elementy iz $\mt$ ne anulirajut se črez $c$, suma jest direktna:
\begin{equation}
        \mo=\mt+\mr.                                             \tag{1}
\end{equation}

Ako teper po uslovju minimalnosti predstavimo $\mo$ kako direktnu sumu direktno nerazložimyh idealov:
\[
        \mo=\mr_1+\cdots+\mr_n,
\]
togda $\mr_i$ morajut byti proste (= minimalne). Bo ako, na primjer, $\mr_1$ ne byl by prosty, a $\mt$ byl by minimalny ideal v $\mr_1$, togda iz (1) imajemo
\[
        \mr_1=\mt+\mr\cap\mr_1 \qquad \hbox{\(\S 4,3\)};
\]
zato $\mr_1$ byl by razložimy, čto ne može byti. Zato $\mo$ jest popolno reducibilno.

Drugo: \emph{kolco bez radikala s uslovjem minimalnosti imaje jediničny element.}

Da najprvo pokažemo suščestvovanje lěvoj jedinice, dostatočno jest pokazati slědujuče: \emph{ako pravy ideal $\ma=c_1\mo\ne\mo$ imaje lěvu jedinicu $c_1$, togda suščestvuje pravy ideal $c\mo$ večej dolgosti, ktory takože imaje lěvu jedinicu $c$.} Bo popolnu reducibilnost, to jest ograničenost vseh dolgostij, jesmo uže pokazali; tak v konečno mnogih krokah dojdemo do lěvoj jedinice za samo $\mo$.

Nehaj zato $\ma=c_1\mo$ i $c_1^2=c_1$. Formula
\[
        r=c_1r+(r-c_1r)
\]
dava, kako vyšje, Peirceovo razloženje $\mo=c_1\mo+\mr$. Nehaj, kako vyšje, $c_2$ bude idempotentny element iz $\mr$, to jest $c_2^2=c_2$, $c_1c_2=0$ (bo $c_1$ anuliraje vse elementy iz $\mr$). Ako položimo $e_1=c_1$, $e_2=c_2-c_2c_1$, togda dostajemo
\[
        e_1^2=e_1,
        \quad e_1e_2=0,
        \quad e_2e_1=0,
        \quad e_2c_2=e_2,
        \quad e_2\ne0,
        \quad e_2^2=e_2.
\]
Po zamětce učinjenoj v primětce, $c=e_1+e_2$ stava lěva jedinica za kolco
\[
        c\mo=e_1\mo+e_2\mo,
\]
koje, zaradi $e_2^2=e_2\ne0$, kako pravy ideal imaje veču dolgost než $\ma$.

Da pokažemo, že konstruovana lěva jedinica $e$ jest takože prava jedinica, činimo Peirceovo razloženje v lěve idealy:
\[
        \mo=\mo e+\ml.
\]
Imajemo $\ml e=0$, $\ml=e\ml$, zato $\ml^2=\ml e\ml=0$, zato, poneže po predpoloženju ne može suščestvovati nilpotentny lěvy ideal, $\ml=0$. Zato $e$ jest takože prava jedinica i tym samym jedinica obće.

\emph{Teorem. Iz pravostrannoj popolnoj reducibilnosti i suščestvovanja jedinice slijedi dvustranna:}
\[
        \mo=\md_1+\cdots+\md_s,
        \qquad \md_i \hbox{ dvustranno proste.}
\]

Kako v prethodnom dokazu, dostatočno jest pokazati, že vsaky minimalny (dvustranny) ideal $\ma$ jest direktny sumand. $\ma$ jest kako pravy ideal direktny sumand:
\[
        \mo=\ma+\mr=e_1\mo+e_2\mo.
\]
Odgovarajuče lěvo razloženje jest
\[
        \mo=\mc+\ml=\mo e_1+\mo e_2.
\]
Slijedi
\[
        \mr\mc\subseteq \mr\ma\leq \mr\cap\ma=0,
        \qquad
        \ml\ma=\mo e_2e_1\mo=0,
\]
\[
        (\ma\ml)^2=\ma\ml\ma\ml=0,
        \qquad \hbox{zato } \ma\ml=0;
\]
\[
        \mr=\mr\mo=(\mr\mc,\mr\ml)=\mr\ml,
        \qquad
        \ml=\mo\ml=(\ma\ml,\mr\ml)=\mr\ml.
\]
Zato $\mr=\ml$, zato $\mr$ jest dvustranny; zato $\ma$ jest dvustranny direktny sumand. Od togo města dokaz ide kako dokaz jednostrannoj popolnoj reducibilnosti.

Idealy $\md_i$ sut takože kako kolca dvustranno proste, bo vse idealy v $\md_i$ sut idealy v $\mo$. Izslědujemo njihovu strukturu.

\end{document}
